\section{Flint--Hills series and its generalization}\label{sec:series}

We first prove Corollary~\ref{thm:flinthills}. Alekseyev showed
that convergence of the classical Flint--Hills series requires
\(\mu(\pi)\le5/2\) \cite[Corollary~4]{alekseyev}; Meiburg proved
that \(\mu(\pi)<5/2\) is sufficient \cite[Theorem~2.5]{meiburg}.
We include the short spacing argument needed for the consequence of
Theorem~\ref{thm:main}. Write
\(\|x\|=\min_{p\in\ZZ}|x-p|\) for distance to the nearest integer.

\begin{lemma}[Dyadic spacing estimate]\label{series:dyadic}
Let \(\alpha>0\) be irrational. Suppose that, for some \(c>0\) and
\(1<\nu<5/2\),
\[
 \|q\alpha\|\ge c q^{1-\nu}\qquad(q\in\ZZ,\ q\ge1).
\]
Then
\[
 \sum_{q=1}^{\infty}\frac{1}{q^3\|q\alpha\|^2}<\infty.
\]
\end{lemma}

\begin{proof}
Fix a positive integer \(K\), and put
\(d=c(2K)^{1-\nu}\). For \(K\le q<2K\), the points
\(q\alpha\) on the circle \(\RR/\ZZ\) have distance at least
\(d\) from zero. Two distinct such points also have circle distance
at least \(d\): apply the hypothesis to their nonzero integer
difference, whose absolute value is less than \(2K\), and use
\(1-\nu<0\).

On either half-circle starting at zero, arrange the distances in
increasing order. Their first distance is at least \(d\), and each
successive distance increases by at least \(d\). Thus these distances
are at least \(d,2d,3d,\ldots\). No positive integer multiple of
\(\alpha\) lies at either endpoint of a half-circle. Consequently
\[
 \sum_{K\le q<2K}\frac{1}{q^3\|q\alpha\|^2}
 \le \frac{2}{K^3d^2}\sum_{j=1}^{\infty}\frac{1}{j^2}
 =O\bigl(K^{2\nu-5}\bigr).
\]
The implied constant depends only on \(c\) and \(\nu\).
Since \(2\nu-5<0\), summing over \(K=1,2,4,\ldots\) proves
the assertion.
\end{proof}

\begin{proof}[Proof of Corollary~\ref{thm:flinthills}]
Fix a real number \(\nu\) with \(2<\nu<5/2\), and choose a
positive integer \(Q\) as in Theorem~\ref{thm:main}. Taking a nearest
integer numerator in that theorem gives
\[
 \|q\pi\|\ge q^{1-\nu}\qquad(q\ge Q).
\]
The same theorem has established irrationality of \(\pi\). Hence
\[
 c=\min\left(\{1\}\cup
       \left\{q^{\nu-1}\|q\pi\|:q\in\ZZ,\ 1\le q<Q\right\}\right)>0.
\]
It follows that \(\|q\pi\|\ge c q^{1-\nu}\) for every positive
integer \(q\), including the finitely many exceptions below \(Q\).
Lemma~\ref{series:dyadic} therefore gives
\begin{equation}\label{series:distance-sum}
 \sum_{q=1}^{\infty}\frac{1}{q^3\|q\pi\|^2}<\infty.
\end{equation}

For each positive integer \(n\), choose a nearest nonnegative integer
\(q\) to \(n/\pi\), and group the integers \(n\) by \(q\).
Then \(|n-q\pi|\le\pi/2\). The group with
\(q=0\) is finite, and its terms in the Flint--Hills series are finite
because irrationality of \(\pi\) excludes \(\sin n=0\) for a
positive integer \(n\). For each \(q\ge1\), its group lies in an
interval of length \(\pi<4\), so contains at most four integers.
Every integer in that group satisfies
\[
 n\ge\pi(q-1/2)\ge\frac{\pi q}{2}.
\]
Concavity of sine on \([0,\pi/2]\) also gives
\[
 |\sin n|=|\sin(n-q\pi)|
 \ge\frac{2}{\pi}|n-q\pi|
 \ge\frac{2}{\pi}\|q\pi\|.
\]
The contribution of the group indexed by \(q\ge1\) is consequently
at most
\[
 \frac{8}{\pi}\frac{1}{q^3\|q\pi\|^2}.
\]
Summing this bound and using \eqref{series:distance-sum} proves
convergence of \(\sum_{n\ge1}n^{-3}\sin^{-2}n\), with angles in
radians.
\end{proof}


The same spacing argument gives the convergence criterion for the
two-parameter family.

\begin{corollary}[Generalized Flint--Hills series]\label{series:generalized}
For fixed real numbers \(a,b>0\), the series
\[
 \sum_{n=1}^{\infty}\frac{1}{n^a|\sin n|^b}
\]
with angles in radians converges if and only if
\(a>\max\{1,b\}\).
\end{corollary}

\begin{proof}
Suppose first that \(a>\max\{1,b\}\), and choose \(\varepsilon>0\)
so that \(a>\max\{1,b\}+b\varepsilon\).
Theorem~\ref{thm:main}, with the finitely many exceptional denominators
absorbed into a positive constant as above, gives
\(\|q\pi\|\ge c_\varepsilon q^{-1-\varepsilon}\) for all \(q\ge1\).
For an integer \(K\ge1\), put \(d=c_\varepsilon(2K)^{-1-\varepsilon}\).
As in the proof of Lemma~\ref{series:dyadic}, the points \(q\pi\)
on \(\RR/\ZZ\) with \(K\le q<2K\) are mutually \(d\)-separated
and at least \(d\) from zero. There are at most \(K\) points, so the half-circle ordering gives
\[
 \sum_{K\le q<2K}\frac{1}{q^a\|q\pi\|^b}
 \le 2K^{-a}d^{-b}\sum_{j=1}^{K}j^{-b}
 \ll_{a,b,\varepsilon}
 \begin{cases}
 K^{1-a+b\varepsilon},&0<b<1,\\
 K^{1-a+\varepsilon}\log(2K),&b=1,\\
 K^{b-a+b\varepsilon},&b>1.
 \end{cases}
\]
Each power of \(K\) is negative by the choice of \(\varepsilon\);
summing over \(K=1,2,4,\ldots\), including the logarithmic factor
when \(b=1\), proves convergence of
\(\sum_{q\ge1}q^{-a}\|q\pi\|^{-b}\).
Use the same nearest-multiple groups as in the preceding proof. The
\(q=0\) group is finite, while the group indexed by \(q\ge1\)
contributes at most
\[
 4(2/\pi)^a(\pi/2)^b q^{-a}\|q\pi\|^{-b}.
\]
This proves the required convergence.

If \(a\le1\), including the boundary \(a=1\), then
\(|\sin n|\le1\) compares the series from below with the divergent
series \(\sum n^{-a}\). It remains to consider \(1<a\le b\).
Let \(p_j/q_j\) be the continued-fraction convergents to \(\pi\),
whose positive integer numerators tend to infinity. They satisfy
\[
 |\sin p_j|=|\sin(p_j-\pi q_j)|
 \le|p_j-\pi q_j|<q_j^{-1}.
\]
Since \(p_j/q_j\to\pi\), eventually \(p_j\le(\pi+1)q_j\), and hence
\[
 \frac{1}{p_j^a|\sin p_j|^b}
 >p_j^{-a}q_j^b\ge(\pi+1)^{-b}p_j^{b-a}.
\]
These summands are bounded away from zero when \(a=b\), and tend
to infinity when \(a<b\). Thus the series also diverges throughout
\(1<a\le b\), including its boundary \(a=b>1\).
\end{proof}
